\documentclass{homework}
% \usepackage{graphicx} % Required for inserting images
\usepackage{amsmath}
\usepackage{gensymb}
\usepackage{subfig}
\usepackage{float}

\class{ECEN 5417}
\title{Project}
\author{Sean Jennings}
\date{December 9, 2023}

\begin{document}
\maketitle

\section{Problem 1}
\begin{letterenum}
\item The solution of the power flow (including both flow and bus data), 
with reactive power enforced is shown in figure \ref{fig:with-limits}:
\begin{center}
    \includegraphics[width=0.8\linewidth]{pf_with_limits.png}
    \captionof{figure}{Power flow with limits}
    \label{fig:with-limits}
\end{center}


\item The solution without limits is shown in figure \ref{fig:without-limits}.
We notice that the powerflow solution without limits enforced has a bus which
exceeds the reactive power limits. Indeed, the generator at bus 8 has a limit of 
24 MVar, which is exceeded by ~0.538 MVar when limits are not enforced. With
limits, the reactive power output by the other generators (at bus 1, 2, 3, and 6)
is increased to compensate.

We also note that the reactive power flowing through each line has also changed.

\item The bus results with increased load is shown in figure \ref{fig:higher-load}.
With the higher power demands, the reactive power limits of every PV bus is
now reached, and as a result, the slack bus must deliver 73MVar to the system, rather
than absorbing reactive power from the system's synchronous condensers.

\item A shunt capacitor of ~0.235p.u. at Bus 14 makes the voltage at bus 14 approximately
1 p.u. This is a vast improvement since without the capacitor, a majority of the nodes have
a voltage below 0.95 p.u., reaching a minimum of 0.899p.u. at bus 14. With the capacitor,
the voltages are now generally much closer to 1p.u., with bus 3 now assuming the
lowest voltage in the system at 0.95p.u.

\item A shunt capacitor of ~0.176p.u. at Bus 13 makes the voltage at bus 13 ~1p.u.
This placement is less effective than the placement at Bus 14.  13 ~1p.u.
Since in the capacitor-less case, bus 14 has the lowest voltage, placing it here allows
the larger sized capacitor to be placed, adding more reactive power support to the
system.

With the placement at Bus 13, we have bus 14 still at a low voltage: 0.95p.u. and
buses 9 and 10 are also relatively low, at 0.96p.u.

The low voltage (and thus the superiority of the placement) at Bus 14 can also
be explained by noting that bus 9 has one of the largest reactive power loads
on the system, at 23.24 MVar.
Additionally, bus 14 is far away from any other generators/condensors.




\begin{center}
    \includegraphics[width=0.8\linewidth]{pf_without_limits.png}
    \captionof{figure}{Power flow without limits}
    \label{fig:without-limits}
\end{center}

\begin{center}
    \includegraphics[width=0.8\linewidth]{higher_load.png}
    \captionof{figure}{Power flow with 40\% higher load}
    \label{fig:higher-load}
\end{center}
\end{letterenum}

\section{Problem 2}

\begin{letterenum}
\item The postive sequence admittance matrix is calculated by taking the "reference"
$Y_{bus}$ matrix and accounting for the self-admittance which is added by loads
and generators. Loads are assumed to be constant impedance so that $Y_{L,k}$, the admittance
due to the load at bus $k$, is calculated as follows:
\begin{equation*}
    Y_{L,k} = \frac{S_k^*}{|V_k|^2}
\end{equation*}
where $S_k\in\C$ is the complex power of the load and $|V_k|\in \R$ is the magnitude of
voltage, also at bus $k$.

We assume that all components has equivalent positive and negative sequence admittances
so that $Y_{bus}^{(+)} = Y_{bus}^{(-)}$: the positive and negative sequence admittance 
matrices are also equal.

The zero-sequence admittance matrix is calculated from line data. Zero sequence current
does not flow through delta-delta or Yg-delta transformers, so only a self-admittance is added
to the Yg side of those transformers. We assume that loads and the Yg-grounded generator are
all well-grounded so that their zero-sequence admittance is equal to the positive sequence.
Delta-connected generators also do not yield 0-sequence admittance.

Final matrices are shown in figure \ref{fig:admittance-matrices}

\begin{center}
    \includegraphics[width=0.8\linewidth]{admittance_matrices.png}
    \captionof{figure}{Positive and Negative Sequence admittance matrices}
    \label{fig:admittance-matrices}
\end{center}

\item Short-circuit fault currents and voltages at each bus are displayed in figure
    \ref{fig:faults}. The 3-phase fault is more severe. As a result of the lower overall
    impedance to ground (all 3 phases have 0 impedance rather than just 1), the fault
    currents are significantly higher.

\begin{center}
    \includegraphics[width=0.8\linewidth]{fault_voltages_and_currents.png}
    \captionof{figure}{Fault voltages and currents for 3-phase and line-to-ground faults}
    \label{fig:faults}
\end{center}

\item Currents in lines (9-14) and (13-14) are shown in Figure \ref{fig:fault_currents}.
The circuit breaker in line 9-14 should have a symmetric interrupting current
less than 3000A, say 2500A, and a single phase interrupting current of less than 1250A.
In line 13-14, the break's symmetric interrupting current should be no more than
2200A and a single phase interrupting current no more than 950A.


\begin{center}
    \includegraphics[width=0.8\linewidth]{fault_currents_in_line.png}
    \captionof{figure}{Fault currents in lines 14-9 and 14-13 for both faults}
    \label{fig:fault_currents}
\end{center}

\end{letterenum}

\section{Problem 3}
\begin{letterenum}
\item Results from running the continuation power flow for
both parts a. and b. are shown in Figure \ref{fig:cpf}.
\item 
\item The continuation power flow makes the case for adding the capacitor
even stronger. From Figure \ref{fig:cpf}, we can see that when line
3-4 goes offline, the curves (both with and without the capacitor)
move to the left. This means, for example, that without the capacitor
any loads between ~4.5-5 MW can be met with the line, but if the line
goes off, the powerflow will fail to converge and the system will
go unstable.

On the other hand, with the capacitor, the curves have moved to the
right, meaning higher loads can be sustained both with and without
the line. In fact, the system with the capacitor and line 3-4 offline
can handle almost as much load at bus 14 as the system without the
capacitor and {\it with} line 3-4!


\begin{center}
    \includegraphics[width=0.8\linewidth]{continuation_power_flow.png}
    \captionof{figure}{Continuation Power flow for parts a. and b.}
    \label{fig:cpf}
\end{center}

\end{letterenum}




\end{document}
